%--- iliad ---
% cluster: B
% title: Training Dynamics
% summary: Exact learning dynamics of deep linear networks -- loss-landscape geometry, balanced gradient flow and the NTK, the rich (saddle-to-saddle) and lazy regimes, their mixed unification, and the implicit bias of SGD noise.
%--- end ---
\documentclass[11pt]{article}

% iliad.sty: local copy first (Overleaf), else the shared ../iliad.sty.
\IfFileExists{iliad.sty}{\usepackage[boxes]{iliad}}{\usepackage[boxes]{../iliad}}

% ---- Free zone -------------------------------------------------------------
% tightlist is emitted by the pandoc source lists we copied verbatim.
\providecommand{\tightlist}{%
  \setlength{\itemsep}{0pt}\setlength{\parskip}{0pt}}
% Unicode glyphs that appear verbatim in the copied prose/table.
\usepackage{newunicodechar}
\newunicodechar{→}{\ensuremath{\rightarrow}}

\title{Training Dynamics}
\author{\authorname{Guillaume Corlouer}\\ \affiliation{Stormglass}}
\date{\today}

\begin{document}
\maketitle

\begin{summary}
Exact learning dynamics of deep linear networks: the geometry of the loss
landscape, balanced gradient flow and the neural tangent kernel, the rich
(saddle-to-saddle) and lazy regimes, their unification via mixed dynamics,
and the implicit bias of SGD noise.
\end{summary}

\begin{learningoutcomes}
  \subsection*{Motivation}
  \begin{itemize}
    \item Know about some big open questions in learning dynamics
    \item Understand the concept of implicit regularization
    \item Know about the different approaches to studying learning dynamics
    \item Understand the AI safety motivations for learning dynamics
    \item Know about emergent misalignment as a safety-relevant phenomena that illustrates the importance of understanding generalization for AI safety
  \end{itemize}

  \subsection*{Geometry and dynamics of deep neural networks}
  \begin{itemize}
    \item Know key results about the loss landscape of deep linear networks (DLNs): critical points are saddles or global minima
    \item Explain the edge of stability phenomenon
    \item Understand that gradient flow can be written as NTK-weighted gradient in function space
    \item DLNs are degenerate and have conserved quantities through gradient flow
  \end{itemize}

  \subsection*{Regimes of learning in deep linear networks}
  \begin{itemize}
    \item Understand the role of initialization, width and depth for the lazy and rich (saddle to saddle) regimes in DLNs
    \item Know about implicit regularization from non-linearity (neural race reduction)
  \end{itemize}
\end{learningoutcomes}

\textbf{References}: \cite{saxe2014}, \cite{achour2024}, \cite{tu2024}.

\section{Setup and notation}

We study deep linear networks (DLNs) with \(L\) weight matrices
\(W_1 \in \mathbb{R}^{d_1 \times d_0}, W_2 \in \mathbb{R}^{d_2 \times d_1}, \ldots, W_L \in \mathbb{R}^{d_L \times d_{L-1}}\).
The network computes:

\[f(x) = W_L W_{L-1} \cdots W_1 \, x =: W x\]

where \(W = W_L \cdots W_1 \in \mathbb{R}^{d_L \times d_0}\) is the
end-to-end (or ``student'') matrix. We train on a dataset
\(\{(x_\mu, y_\mu)\}_{\mu=1}^N\) with the squared loss:

\[\mathcal{L}(\theta) = \frac{1}{2N} \sum_{\mu=1}^{N} \|y_\mu - W_L \cdots W_1 \, x_\mu\|^2\]

In the population limit with whitened inputs \(\Sigma_X = I\), this
becomes:

\[\mathcal{L}(W) = \frac{1}{2}\|M - W\|_F^2\]

where \(M = \Sigma_{YX} \Sigma_X^{-1}\) is the ``teacher'' matrix (the
OLS solution) with SVD
\(M = U \,\text{diag}(s_1, \ldots, s_r, 0, \ldots, 0)\, V^\top\), and
\(s_1 \geq s_2 \geq \cdots \geq s_r > 0\).

\section{Loss landscape geometry}
\label{sec:td-geometry}

This problem explores the critical point structure of deep linear
networks, following \cite{achour2024}.

\begin{exercise}[Diagonal decomposition]
\label{ex:td-1a}
Consider the \textbf{diagonal} case:
\(d_0 = d_1 = \cdots = d_L = d\), and the teacher is diagonal,
\(M = \text{diag}(s_1, \ldots, s_d)\), with
\(s_1 > s_2 > \cdots > s_d > 0\). Restrict attention to diagonal weight
matrices \(W_l = \text{diag}(w_l^{(1)}, \ldots, w_l^{(d)})\). Show that
the loss decomposes into \(d\) independent scalar problems:

\[\mathcal{L} = \frac{1}{2}\sum_{\alpha=1}^d \left(s_\alpha - \prod_{l=1}^L w_l^{(\alpha)}\right)^2\]
\end{exercise}


\begin{exercise}[Scalar critical points]
\label{ex:td-1b}
For a single scalar mode with target \(s > 0\), find all
first-order critical points of
\(\ell(w_1, w_2) = \frac{1}{2}(s - w_1 w_2)^2\) at depth \(L = 2\). Show
that the critical points are:

\begin{itemize}
\tightlist
\item
  The global minimum manifold: \(w_1 w_2 = s\)
\item
  The origin: \(w_1 = w_2 = 0\)
\end{itemize}

Classify the origin as a saddle point by computing the Hessian \(H\) of
\(\ell\) at \((0,0)\) and showing it has both positive and negative
eigenvalues.

\emph{Hint}: Write the gradient equations
\end{exercise}


\begin{exercise}[Critical-point structure]
\label{ex:td-1c}
Achour et al.~\cite{achour2024} show that every
first-order critical point \(\theta = (W_1, \ldots, W_L)\) satisfies:
there exists a subset \(S \subseteq \{1, \ldots, d\}\) such that

\[W = W_L \cdots W_1 = P_S\, M\]

where \(P_S = U_S U_S^\top\) is the orthogonal projector onto the span
of the left singular vectors \(\{u_\alpha\}_{\alpha \in S}\) of \(M\).
For the diagonal case \(M = \text{diag}(s_1, \ldots, s_d)\), write the
value of the loss at the critical point corresponding to
\(S = \{1, \ldots, r\}\) with \(r < d\).
\end{exercise}


\begin{exercise}[Symmetry of the global minima]
\label{ex:td-1d}
The group
\(GL_h := GL_{d_1} \times \cdots \times GL_{d_{L-1}}\) acts on the
weights by:

\[(W_1, \ldots, W_L) \mapsto (g_1 W_1,\; g_2 W_2 g_1^{-1},\; \ldots,\; W_L g_{L-1}^{-1})\]

Verify that the student map \(\mu(\theta) = W_L \cdots W_1\) is
invariant under this action: \(\mu(g \cdot \theta) = \mu(\theta)\). What
does this imply about the dimension of the set of global minima in
parameter space?
\end{exercise}


\section{Gradient flow and conserved quantities}
\label{sec:td-gradflow}

We now study gradient flow:
\(\dot{\theta}(t) = -\nabla_\theta \mathcal{L}(\theta(t))\), the
continuous-time limit of gradient descent with infinitesimal learning
rate.

\begin{exercise}[Gradient-flow equations]
\label{ex:td-2a}
For a two-layer diagonal DLN (\(L=2\)) with loss
\(\mathcal{L} = \frac{1}{2}\|M - W_2 W_1\|_F^2\), assume that each weight matrix $W_i$ is a diagonal matrix and whitened inputs.
Derive the gradient flow equations for each layer. Show that:

\[\dot{W}_1 = W_2^\top(M - W_2 W_1), \qquad \dot{W}_2 = (M - W_2 W_1) W_1^\top\]
\end{exercise}


\begin{exercise}[Balancedness is conserved]
\label{ex:td-2b}
Define the \textbf{balancedness matrix}:

\[G := W_2^\top W_2 - W_1 W_1^\top\]

Show that \(G\) is conserved under the gradient flow,
i.e.~\(\dot{G} = 0\). In other words, gradient flow is constrained on the balanced manifold \begin{equation*}
    \mathcal{M}_{\beta}
    :=
    \{\theta\in\Omega\mid
    W_{\ell+1}^{\top}W_{\ell+1}
    -
    W_{\ell}W_{\ell}^{\top}
    =
    \beta_\ell,\;
    \ell=1,\ldots,L-1\}.
\end{equation*}

\emph{Hint}: Compute
\(\dot{G}\),
substitute the gradient flow equations and verify that the
terms cancel pairwise.
\end{exercise}


From now on, assume \textbf{balanced initialization}:
\(G(0) = 0\), which by \cref{ex:td-2b} means \(W_2^\top W_2 = W_1 W_1^\top\)
for all time. We want to derive the gradient flow in function space (the
ODE for the student \(W = W_2 W_1\)). This requires several steps.

\begin{exercise}[Function-space velocity]
\label{ex:td-2ci}
Compute \(\dot{W} := \frac{d}{dt}(W_2 W_1)\) using the
gradient flow equations from \cref{ex:td-2a}. Show that:

\[\dot{W} = (M - W)\, W_1^\top W_1 + W_2 W_2^\top\, (M - W)\]
\end{exercise}


\begin{exercise}[\(W_1^\top W_1\) in terms of \(W\)]
\label{ex:td-2cii}
We now need to express \(W_1^\top W_1\) and
\(W_2 W_2^\top\) in terms of \(W = W_2 W_1\). For simplicity, assume
that the weight matrices \(W_l\) are diagonal. Show that:

\[W_1^\top W_1 = (W^\top W)^{1/2}\]
\end{exercise}


\begin{exercise}[\(W_2 W_2^\top\) in terms of \(W\)]
\label{ex:td-2ciii}
Similarly, show that
\(W_2 W_2^\top = (W W^\top)^{1/2}\).
\end{exercise}


\begin{exercise}[The balanced function-space ODE]
\label{ex:td-2civ}
Substitute the results of \cref{ex:td-2cii,ex:td-2ciii} into \cref{ex:td-2ci}
to obtain:

\[\dot{W} = (W W^\top)^{1/2}(M - W) + (M - W)(W^\top W)^{1/2}\]
\end{exercise}


\begin{exercise}[The NTK operator]
\label{ex:td-2d}
Define the NTK operator for \(L = 2\) as:

\[K[F] := (WW^\top)^{1/2}\, F + F\, (W^\top W)^{1/2}\]

Show that the gradient flow from \cref{ex:td-2civ} can be written as
\(\dot{W} = K[M - W]\).

It turns out that this result generalizes to depth \(L\) on the balanced
manifold and to non-diagonal weight matrices:

\[\dot{W} = \sum_{k=1}^{L} (WW^\top)^{\frac{L-k}{L}} (M - W) (W^\top W)^{\frac{k-1}{L}}\]
This is the NTK equation in the case of DLNs. The NTK equation is a gradient flow in function space with NTK being a preconditioning operator for the gradient.
\end{exercise}


\section{Rich regime: saddle-to-saddle}
\label{sec:td-rich}

This is the core problem. We derive the exact solution of the rich
regime dynamics directly from the self-consistent equation of \cref{sec:td-gradflow},
emphasizing the NTK perspective.

\textbf{Alignment assumption.} Start from the balanced
gradient flow equation derived in \cref{ex:td-2civ}:

\[\dot{W} = (WW^\top)^{1/2}(M - W) + (M - W)(W^\top W)^{1/2}\]

We work in the \textbf{rich regime} (small initialization) with balanced
weights. Assume that the student is \textbf{aligned} to the teacher: the
singular vectors of \(W(t)\) coincide with those of \(M\) at all times.
Concretely, write:

\[M = U\,\text{diag}(s_1, \ldots, s_d)\,V^\top, \qquad W(t) = U\,\text{diag}(w_1(t), \ldots, w_d(t))\,V^\top\]

where \(U, V\) are the left/right singular vectors of \(M\), and
\(w_\alpha(t) \geq 0\) are the evolving singular values of \(W\).

\begin{exercise}[Aligned NTK]
\label{ex:td-3ai}
Under this alignment assumption, show that the NTK
operator is aligned to the task, i.e.~\(K[F]\) preserves the SVD basis.
In particular, show that:

\[(WW^\top)^{1/2} = U\,\text{diag}(w_1, \ldots, w_d)\,U^\top, \qquad (W^\top W)^{1/2} = V\,\text{diag}(w_1, \ldots, w_d)\,V^\top\]

and that the residual is
\(M - W = U\,\text{diag}(s_1 - w_1, \ldots, s_d - w_d)\,V^\top\).

\emph{Hint}: Recall that for \(A = U\,\text{diag}(\sigma_i)\,V^\top\),
we have \(AA^\top = U\,\text{diag}(\sigma_i^2)\,U^\top\), and therefore
\((AA^\top)^{1/2} = U\,\text{diag}(|\sigma_i|)\,U^\top\).
\end{exercise}


\begin{exercise}[Decoupled scalar ODEs]
\label{ex:td-3aii}
Substitute into the self-consistent equation. Show that
the matrix equation decouples into \(d\) independent scalar ODEs:

\[\dot{w}_\alpha = 2\,w_\alpha\,(s_\alpha - w_\alpha), \qquad \alpha = 1, \ldots, d\]

\emph{Hint}: Compute the product \((WW^\top)^{1/2}(M - W)\) in the SVD
basis. It is diagonal with entries \(w_\alpha(s_\alpha - w_\alpha)\).
The second term contributes identically, giving the factor of 2.
\end{exercise}


\begin{exercise}[Timescale of learning]
\label{ex:td-3b}
The ODE
\(\dot{w} = 2w(s - w)\) is a logistic equation whose solution is:

\[w(t) = \frac{s}{1 + \left(\frac{s}{w_0} - 1\right)e^{-2st}}\]

where \(w_0 = w(0)\). Compute the time \(t_\alpha\) it takes for mode
\(\alpha\) to travel from initial strength \(w_0\) to a final strength
\(w_f\). Show that:

\[t_\alpha = \frac{1}{2 s_\alpha} \ln\left(\frac{w_f\,(s_\alpha - w_0)}{w_0\,(s_\alpha - w_f)}\right)\]

Deduce that modes with larger singular values \(s_\alpha\) are learned
faster. This is a \textbf{separation of timescales}: the network learns
features in decreasing order of their singular value strength.
\end{exercise}


\begin{exercise}[Incremental learning]
\label{ex:td-3c}
Consider a teacher with
\(r\) nonzero singular values and a small uniform initialization
\(w_\alpha(0) = w_0 \ll s_r\) for all \(\alpha\). Discuss:

\begin{itemize}
\tightlist
\item
  What is the approximate rank of \(W(t)\) at early times?
\item
  How does the rank evolve at intermediate times?
\item
  What happens in the limit \(t \to \infty\)?
\end{itemize}

What does this tell us about the sequence of critical points visited by
the gradient flow?
\end{exercise}


\section{Lazy regime}
\label{sec:td-lazy}

We now show that large initialization freezes the NTK and eliminates the
timescale separation found in \cref{sec:td-rich}. For clarity, we work with the
\textbf{diagonal scalar} model from \cref{ex:td-1a}, so each mode
\(\alpha\) is independent. This avoids matrix algebra and isolates the
essential mechanism.

\textbf{Setup.} Consider a single mode: a depth-2 diagonal
DLN with scalar weights \(a, b\) learning a target \(s > 0\). From
\cref{sec:td-gradflow,sec:td-rich}, the balanced gradient flow for \(w = ab\) is:

\[\dot{w} = 2w(s - w)\]

The factor \(2w\) is the (scalar) NTK --- it is
\textbf{state-dependent}: the effective learning rate depends on the
current value of \(w\).

\begin{exercise}[Linearizing the NTK]
\label{ex:td-4a}
Now initialize at a \textbf{large} value \(w_0 \gg s > 0\). Assume that in
the early phase of training, while \(w\) has not changed much from
\(w_0\), the ODE becomes approximately:

\[\dot{w} \approx 2w_0(s - w)\]
\end{exercise}


\begin{exercise}[Frozen-NTK solution]
\label{ex:td-4b}
Solve the linearized ODE from \cref{ex:td-4a}. Show that:

\[w(t) = s + (w_0 - s)\,e^{-2w_0 t}\]

What is the learning timescale?
\end{exercise}


\begin{exercise}[No timescale separation]
\label{ex:td-4c}
Now restore the mode
index \(\alpha\). In the lazy regime, each mode satisfies
\(\dot{w}_\alpha \approx 2w_0(s_\alpha - w_\alpha)\) with the
\textbf{same} rate \(2w_0\) for all \(\alpha\). Compare the lazy
learning time \(t_\alpha^{\text{lazy}}\) with the rich-regime timescale
from \cref{ex:td-3b}: \(t_\alpha^{\text{rich}} \propto 1/s_\alpha\).

Discuss the difference in \textbf{rank bias} between the lazy and rich
regimes.
\end{exercise}


\section{Mixed dynamics: unifying lazy and rich}
\label{sec:td-mixed}

In practice, networks are neither purely lazy nor purely rich. Tu,
Aranguri \& Jacot (2024) provide a unified description. We develop the
key ideas using the diagonal scalar model.

\begin{exercise}[The interpolating ODE]
\label{ex:td-5a}
In \cref{sec:td-rich,sec:td-lazy}, we
studied the same ODE \(\dot{w}_\alpha = 2w_\alpha(s_\alpha - w_\alpha)\)
in two limits: small \(w_0\) (rich) and large \(w_0\) (lazy). A more
general model for a depth-2 DLN with balanced initialization at scale
\(\sigma\) and width \(n\) replaces the scalar NTK \(2w_\alpha\) by:

\[\dot{w}_\alpha = 2\sqrt{w_\alpha^2 + \tau^2}\;(s_\alpha - w_\alpha)\]

where \(\tau := \sigma^2 \sqrt{n}\) is a \textbf{threshold} parameter
that depends on the initialization scale and width.

Verify that this ODE reproduces the two known regimes:

\begin{itemize}
\tightlist
\item
  \textbf{Rich limit} (\(\tau \to 0\)): recover
  \(\dot{w}_\alpha = 2|w_\alpha|(s_\alpha - w_\alpha)\).
\item
  \textbf{Lazy limit} (\(\tau \to \infty\) with \(w_\alpha\) bounded):
  recover \(\dot{w}_\alpha \approx 2\tau(s_\alpha - w_\alpha)\).
\end{itemize}
\end{exercise}


\begin{exercise}[Two-phase dynamics]
\label{ex:td-5b}
At initialization, all modes
start at \(w_\alpha(0) \sim \sigma^2 \ll \tau\) (since
\(\sqrt{n} \gg 1\)). Observe that:

\begin{itemize}
\tightlist
\item
  When \(|w_\alpha| \ll \tau\), the effective learning rate is
  \(\approx 2\tau\), the same for all modes (lazy behavior).
\item
  When \(|w_\alpha| \gg \tau\), the effective learning rate is
  \(\approx 2|w_\alpha|\), which is mode-dependent (rich behavior).
\end{itemize}

Describe the resulting two-phase dynamics in words: what happens first,
and what happens later?
\end{exercise}


\begin{exercise}[Mode-by-mode transition]
\label{ex:td-5c}
Not all modes cross the
threshold \(\tau\) at the same time: which modes cross it first? Explain
why this means that a single network can simultaneously have some modes
in the rich regime and others still in the lazy regime.
\end{exercise}


\section{Stochastic implicit bias (bonus)}
\label{sec:td-stochastic}

SGD introduces noise from mini-batching. Write the SGD update as a combination of gradient and noise term (use a LLM to check your answer). A continuous model is the
Langevin SDE:

\[d\theta_t = -\nabla \mathcal{L}(\theta_t)\,dt + \sqrt{\eta\,\Sigma(\theta_t)}\,dB_t\]

where \(\Sigma(\theta)\) is the covariance of the stochastic gradient
noise and \(\eta\) is the learning rate.

\begin{exercise}[Boltzmann equilibrium]
\label{ex:td-6a}
The probability density \(p(\theta, t)\) of the parameters
evolves according to the Fokker--Planck equation:

\[\partial_t p = -\nabla \cdot \mathbf{j}, \qquad \mathbf{j} = -\nabla \mathcal{L}(\theta)\,p(\theta) - \frac{\eta}{2}\nabla \cdot \left(\Sigma(\theta)\,p(\theta)\right)\]

where \(\mathbf{j}\) is the probability current. Assume: (i)
stationarity \(\partial_t p^* = 0\), (ii) thermal equilibrium
\(\mathbf{j} = 0\), and (iii) isotropic noise \(\Sigma = \sigma^2 I\).
Show that the equilibrium distribution is the Boltzmann distribution:

\[p^*(\theta) \propto \exp\left(-\frac{2}{\eta \sigma^2}\mathcal{L}(\theta)\right)\]

\emph{Hint}: Setting \(\mathbf{j} = 0\) with \(\Sigma = \sigma^2 I\)
gives \(\nabla \mathcal{L}\, p + \frac{\eta\sigma^2}{2}\nabla p = 0\).
This is a first-order ODE for \(p\) in terms of \(\mathcal{L}\). Try the
ansatz \(p \propto e^{-\beta \mathcal{L}}\) and solve for \(\beta\).
\end{exercise}


\begin{exercise}[Temperature and flatness]
\label{ex:td-6b}
The ratio \(\frac{2}{\eta \sigma^2}\) plays the role of an
inverse temperature \(\beta\). Interpret what happens to the equilibrium
distribution when:

\begin{itemize}
\tightlist
\item
  \(\eta\) is very small (low temperature)
\item
  \(\eta\) is very large (high temperature)
\end{itemize}

Which regime favors flatter minima, and why might this be beneficial for
generalization?
\end{exercise}


\begin{exercise}[Anisotropic noise]
\label{ex:td-6c}
In practice, SGD noise is \textbf{not} isotropic:
\(\Sigma(\theta)\) depends on both the loss landscape and the data.
Without solving anything, explain qualitatively why anisotropic noise
can introduce an implicit bias that goes beyond what the loss function
\(\mathcal{L}\) alone would select. Specifically, why might SGD
preferentially escape sharp directions of the loss while remaining
stable along flat directions?
\end{exercise}


\section*{Learn more}

The readings are roughly ordered in a way that makes sense for learning. This is a curated list of what I find important, and it is not exhaustive.

\subsection*{Key readings}
\begin{itemize}
  \item \href{https://arxiv.org/abs/2107.13289}{The loss landscape of deep linear neural networks: a second-order analysis} Achour et al. --- First-order and second-order classification of critical points in DLN loss landscapes. Classifies strict and non-strict saddles.
  \item Saxe et al. \href{https://www.pnas.org/doi/10.1073/pnas.1820226116}{A mathematical theory of semantic development in deep neural networks} --- Read the supplementary materials section to understand the exact solution of gradient flow dynamics in deep linear networks
  \item \href{https://arxiv.org/abs/2411.19920}{Geometry of fibers of the multiplication map of deep linear neural networks} Simon Pepin Lehalleur et al. --- Stratifies global minima into orbits using quiver representation theory
  \item \href{https://arxiv.org/abs/2411.09004}{The geometry of the deep linear network} Govind Menon --- Beautiful mathematical treatment of gradient flow in DLNs and a surprising mathematical result relating gradient flow and free-energy minimization. Makes implicit regularization explicit (under balanced assumptions)
  \item \href{https://arxiv.org/abs/2602.07852}{Emergent Misalignment is Easy, Narrow Misalignment is Hard} --- When fine-tuned on data with harmful inputs, the model generalize in harmful ways on other unrelated datasets.
  \item \href{https://arxiv.org/abs/2602.07852}{Emergent Misalignment is Easy, Narrow Misalignment is Hard} --- By regularizing, we can mitigate emergent misalignment
  \item Literature Review: \href{https://arxiv.org/abs/2604.21691}{There Will Be a Scientific Theory of Deep Learning}
  \item \href{https://arxiv.org/abs/2209.15594}{Self-Stabilization: The Implicit Bias of Gradient Descent at the Edge of Stability}
  \item \href{https://arxiv.org/abs/2306.08590}{Beyond Implicit Bias: The Insignificance of SGD Noise in Online Learning}
\end{itemize}

\subsection*{Loss landscape geometry}
\begin{itemize}
  \item \href{https://arxiv.org/abs/1605.07110}{Deep Learning without Poor Local Minima}, Kenji Kawaguchi --- For non-bottlenecked DLNs, all local minima are global
  \item \href{https://arxiv.org/abs/2107.13289}{The loss landscape of deep linear neural networks: a second-order analysis} Achour et al. --- First-order and second-order classification of critical points in DLN loss landscapes. Classifies strict and non-strict saddles.
  \item \href{https://arxiv.org/abs/1910.01671}{Pure and Spurious Critical Points: a Geometric Study of Linear Networks} Matthew Trager et al. --- Defines spurious minima and counts the number of connected components of the global minima
  \item \href{https://arxiv.org/abs/2411.19920}{Geometry of fibers of the multiplication map of deep linear neural networks} Simon Pepin Lehalleur et al. --- Stratifies global minima into orbits using quiver representation theory
  \item \href{https://arxiv.org/abs/1412.0233}{The Loss Surfaces of Multilayer Networks} Anna Choromanska et al --- No-local-minima results in non-linear multilayer networks under some strong assumptions (spin glass models)
  \item \href{https://arxiv.org/abs/1802.10026}{Loss Surfaces, Mode Connectivity, and Fast Ensembling of DNNs}, Timur Garipov et al --- Study the connectedness of global minima of loss landscapes (mode connectivity)
  \item \href{https://arxiv.org/abs/1712.10132}{The Multilinear Structure of ReLU Networks}, Thomas Laurent, James von Brecht --- Show that local minima are typically singular
  \item Classic result: \href{https://www.sciencedirect.com/science/article/abs/pii/0893608089900142}{Neural networks and principal component analysis: Learning from examples without local minima} Pierre Baldi et al. --- Original result that linear network landscapes have no spurious local minima (single hidden layer case)
  \item Good review of key DLNs results: \href{https://arxiv.org/abs/2511.10362}{Gradient Flow Equations for Deep Linear Neural Networks: A Survey from a Network Perspective}, Joel Wendin, Claudio Altafini --- Also covers gradient flow
\end{itemize}

\subsection*{Implicit biases of gradient flow}
\begin{itemize}
  \item Saxe et al. \href{https://www.pnas.org/doi/10.1073/pnas.1820226116}{A mathematical theory of semantic development in deep neural networks} --- Read the supplementary materials section to understand the exact solution of gradient flow dynamics in deep linear networks
  \item \href{https://arxiv.org/abs/2106.15933}{Saddle-to-Saddle Dynamics in Deep Linear Networks: Small Initialization Training, Symmetry, and Sparsity} Arthur Jacot et al. --- Studies the saddle-to-saddle dynamics in DLNs. Introduces the different regimes (lazy and rich)
  \item \href{https://arxiv.org/abs/2411.09004}{The geometry of the deep linear network} Govind Menon --- Beautiful mathematical treatment of gradient flow in DLNs and a surprising mathematical result relating gradient flow and free-energy minimization. Makes implicit regularization explicit (under balanced assumptions)
  \item \href{https://arxiv.org/abs/2512.20607}{Saddle-to-Saddle Dynamics Explains A Simplicity Bias Across Neural Network Architectures} Saxe et al. --- Simplicity bias of gradient flow from DLNs extends to non-linear and transformer architectures under small initialization
  \item \href{https://arxiv.org/abs/2307.00144}{Abide by the Law and Follow the Flow: Conservation Laws for Gradient Flows} --- Gives a recipe to exhaustively derive conservation laws through gradient flow
  \item \href{https://arxiv.org/abs/2302.11055}{SGD learning on neural networks: leap complexity and saddle-to-saddle dynamics} --- SGD builds up complex solutions by first learning simpler solutions and then composing them
  \item \href{https://arxiv.org/abs/2405.17580}{Mixed Dynamics In Linear Networks: Unifying the Lazy and Active Regimes} --- Unifies lazy and rich in 2-layer DLNs and gives conditions for the transition between them
  \item \href{https://arxiv.org/abs/1806.07572}{Neural Tangent Kernel: Convergence and Generalization in Neural Networks} Arthur Jacot, Franck Gabriel, Clément Hongler --- NTK paper that describes gradient flow in function space
  \item \href{https://arxiv.org/abs/2207.10430}{The Neural Race Reduction: Dynamics of Abstraction in Gated Networks} Saxe et al. --- Implicit biases toward shared representation in non-linear networks as modelled by gated DLNs
  \item \href{https://arxiv.org/abs/2506.06489}{Alternating Gradient Flows: A Theory of Feature Learning in Two-layer Neural Networks} Daniel Kunin et al. --- A theory of feature learning as a two-step process between dormant and active neurons
  \item \href{https://arxiv.org/abs/2409.14623}{From Lazy to Rich: Exact Learning Dynamics in Deep Linear Networks} --- Studies the transition between lazy and rich regimes in DLNs with balancedness parameters
  \item \href{https://arxiv.org/abs/1705.09280}{Implicit Regularization in Matrix Factorization} --- Gradient descent on matrix factorization converges to the minimum nuclear norm
  \item \href{https://arxiv.org/abs/1906.05890}{Gradient Descent Maximizes the Margin of Homogeneous Neural Networks} --- Implicit bias toward max-margin in classification
  \item \href{https://arxiv.org/abs/1810.02281}{A Convergence Analysis of Gradient Descent for Deep Linear Neural Networks} Sanjeev Arora et al. --- Convergence analysis of SGD
\end{itemize}

\subsection*{Learning rate (discrete GD)}
\begin{itemize}
  \item \href{https://arxiv.org/abs/2209.15594}{Self-Stabilization: The Implicit Bias of Gradient Descent at the Edge of Stability} --- Shows that curvature stabilizes around the inverse of the learning rate.
  \item \href{https://arxiv.org/abs/2410.24206}{Understanding Optimization in Deep Learning with Central Flows} --- GD with a discrete learning rate is equivalent to gradient flow on an effective loss
  \item \href{https://arxiv.org/abs/2410.05192}{Understanding Warmup-Stable-Decay Learning Rates: A River Valley Loss Landscape Perspective} --- Pretraining exhibits a river valley loss landscape and gives intuition about the learning-rate schedule (warm-up, stable, decay)
  \item \href{https://www.nature.com/articles/s41467-025-58532-9}{Optimization on multifractal loss landscapes explains a diverse range of geometrical and dynamical properties of deep learning} --- Models the landscape as multifractal and analyses sub- and super-diffusive behaviour of GD
\end{itemize}

\subsection*{Stochasticity}
\begin{itemize}
  \item \href{https://arxiv.org/abs/2602.12257}{On the implicit regularization of Langevin dynamics with projected noise} --- Makes explicit the implicit bias of stochasticity in deep linear networks with a Langevin model
  \item \href{https://arxiv.org/abs/2306.08590}{Beyond Implicit Bias: The Insignificance of SGD Noise in Online Learning} --- Golden Path Hypothesis: in the transient regime dominated by drift (typical of pretraining), SGD is simply a noisy deformation of GD (it does not select different basins)
  \item \href{https://arxiv.org/abs/2306.04251}{Stochastic Collapse: How Gradient Noise Attracts SGD Dynamics Towards Simpler Subnetworks} --- Stochasticity induces a bias toward simpler solutions in deep linear networks
  \item \href{https://arxiv.org/abs/2109.14119}{Stochastic Training is Not Necessary for Generalization} --- Makes the implicit regularization of stochasticity explicit
  \item \href{https://arxiv.org/abs/2106.09524}{Implicit Bias of SGD for Diagonal Linear Networks: a Provable Benefit of Stochasticity} --- Shows that stochasticity has a bias toward sparser solutions relative to GD
  \item \href{https://arxiv.org/abs/1710.11029}{Stochastic gradient descent performs variational inference, converges to limit cycles for deep networks} --- SGD minimizes some free energy and can have circular currents at convergence
  \item \href{https://arxiv.org/abs/1704.04289}{Stochastic Gradient Descent as Approximate Bayesian Inference} --- Classic paper on SGD locally approximating Bayesian inference, although it assumes non-degeneracies
  \item \href{https://arxiv.org/abs/2002.03495}{A Diffusion Theory For Deep Learning Dynamics: Stochastic Gradient Descent Exponentially Favors Flat Minima} --- Implicit bias toward flat minima
  \item \href{https://arxiv.org/abs/2205.07069}{Implicit Regularization or Implicit Conditioning? Exact Risk Trajectories of SGD in High Dimensions} --- Shows the Golden Path Hypothesis on quadratic loss in a convex setup using an interesting SDE model of SGD.
  \item \href{https://arxiv.org/abs/1812.06162}{An Empirical Model of Large-Batch Training} --- Gradient noise scale can be used to decide the optimal batch size
  \item \href{https://arxiv.org/abs/2503.22478}{Almost Bayesian: The Fractal Dynamics of Stochastic Gradient Descent} --- Relates SGD to Bayesian training
  \item \href{https://arxiv.org/abs/2306.04815}{Catapults in SGD: spikes in the training loss and their impact on generalization through feature learning} --- Loss spikes when training with SGD impact generalization
  \item \href{https://arxiv.org/abs/2006.04740}{The Heavy-Tail Phenomenon in SGD} --- SGD noise can be heavy-tailed, and this induces a different regularization
\end{itemize}

\subsection*{Emergence: empirical examples}
\begin{itemize}
  \item \href{https://arxiv.org/abs/2201.02177}{Grokking: Generalization Beyond Overfitting on Small Algorithmic Datasets} --- Grokking: delayed generalization on modular addition
  \item \href{https://transformer-circuits.pub/2022/in-context-learning-and-induction-heads/index.html}{In-context Learning and Induction Heads} --- Induction heads form during training. They are important for in-context learning
  \item \href{https://arxiv.org/abs/2502.17424}{Emergent Misalignment: Narrow finetuning can produce broadly misaligned LLMs} --- Emergent misalignment: fine-tuning on insecure code can generalize to misaligned behaviour on other data
  \item \href{https://arxiv.org/abs/2602.07852}{Emergent Misalignment is Easy, Narrow Misalignment is Hard} --- By regularizing, we can mitigate emergent misalignment
  \item \href{https://arxiv.org/abs/2111.00034}{Neural Networks as Kernel Learners: The Silent Alignment Effect} --- While the loss is flat, student singular vectors align with teacher singular vectors, and this can be detected with the NTK
  \item \href{https://arxiv.org/abs/2304.15004}{Are Emergent Abilities of Large Language Models a Mirage?} --- Emergence is in the eye of the metric used to measure it
  \item \href{https://arxiv.org/abs/2203.15556}{Training Compute-Optimal Large Language Models} --- Chinchilla scaling law paper
  \item \href{https://arxiv.org/abs/2206.07682}{Emergent Abilities of Large Language Models} --- Shows that novel capabilities emerge with more compute
\end{itemize}

\subsection*{Theoretical approaches to emergence}
\begin{itemize}
  \item \href{https://arxiv.org/abs/2310.06110}{Grokking as the Transition from Lazy to Rich Training Dynamics}
  \item \href{https://arxiv.org/abs/2310.03789}{Grokking as a First Order Phase Transition in Two Layer Networks, Rubin et al.}
  \item \href{https://www.youtube.com/watch?v=WcWFFiPRslM&t=1056s}{Blake Bordelon - Infinite limits and scaling laws of neural networks - IPAM at UCLA}
  \item \href{https://arxiv.org/abs/2307.15936}{A Theory for Emergence of Complex Skills in Language Models}
  \item \href{https://ericjmichaud.com/quanta/}{On neural scaling and the quanta hypothesis}
  \item \href{https://pehlevan.seas.harvard.edu/sites/g/files/omnuum6471/files/pehlevan/files/princeton_lecture_notes.pdf}{Lecture Notes on Infinite-Width Limits of Neural Networks; Cengiz Pehlevan and Blake Bordelon}
  \item \href{https://arxiv.org/abs/2601.01010}{Disordered Dynamics in High Dimensions: Connections to Random Matrices and Machine Learning; Blake Bordelon, Cengiz Pehlevan}
  \item \href{https://arxiv.org/abs/2502.18553}{Applications of Statistical Field Theory in Deep Learning; Zohar Ringel et al.}
  \item \href{https://link.springer.com/book/10.1007/978-3-030-46444-8}{Statistical Field Theory for Neural Networks, Moritz Helias, David Dahmen}
  \item \href{https://arxiv.org/abs/2602.12855}{Lecture notes: From Gaussian processes to feature learning, Moritz Helias et al}
\end{itemize}


\bibliographystyle{alphaurl}
\bibliography{biblo}

\end{document}
